% New writing for the introduction of Volume I. This file is never overwritten by split_volumes.py.

\section{The two volumes}\label{ov1:sec:volumes}

This volume is the theory, and its companion, Volume II, is the computation. Part I gives every representation of $\Xi$ the two volumes use: the analytic ones in Chapter~\ref{vol1:ch2}, the probabilistic ones in Chapter~\ref{vol1:ch3}, the classes of infinitely divisible laws in which they live (generalised gamma convolutions, their extended form on the line, the Bondesson hierarchy and hyperbolic complete monotonicity) in Chapter~\ref{vol1:ch4}, and the reciprocal $\Xi$ function, whose Thorin measure is the subject of both volumes, in Chapter~\ref{vol1:ch5}. Parts II to V follow the Thorin condition through the reciprocal zeta function, the basepoints and the critical strip, other $L$ functions, and the probabilistic and statistical mechanical readings. Volume II computes the objects defined here, from the P\'olya kernel and in ball arithmetic, and its central part runs the Cauchy flow backwards: the Thorin measure at level $\theta$ is the Cauchy smoothing of the zeros, and the question whether that smoothing can be undone down to the critical line is RH.

References into Volume II are printed with their numbers there.

\section{Conventions}\label{ov1:sec:conventions}

Every statement carries one of six words. \emph{Proved}: a proof or a pointer to one is given. \emph{Computed}: a script of the book produced the number, in floating point at the stated precision. The scripts and their data are not named in the text: they are in the book's repository, \url{https://github.com/VadimSokolov/computing_rh}, where \texttt{code/README.md} lists them by chapter. \emph{Certified}: the number was computed in ball arithmetic and the enclosure is rigorous. \emph{Conditional}: the statement assumes RH, GRH, or the zero free half plane $\re s>\frac78$ announced in October 2026 and not refereed; the hypothesis is named each time. \emph{Heuristic} and \emph{open} mean what they say. No statement of either volume proves RH.

One density appears throughout under the names of the chapters that use it. With $\alpha=\theta=\sigma=\frac12+\varepsilon$,
\[
\rho_\varepsilon(t)=\frac1\pi\,v_\alpha(t),\qquad v_\alpha(t)=\re\frac{\xi'}{\xi}(\alpha+it),\qquad u_\theta(v)=\frac1\pi\re\frac{\xi'}{\xi}(\theta+iv)=m_\theta(v),
\]
so that $\rho_\varepsilon$ (the Thorin density at resolution $\varepsilon$), $v_\alpha/\pi$ (the density at the basepoint $\alpha$) and $u_\theta$ (the density at level $\theta$) are one function in three parametrisations. This volume writes basepoints $\alpha$ and the bracket $v_\alpha$; Volume II writes levels $\theta$ and the density $u_\theta$. We write $b_\alpha=v_\alpha(0)=\frac{\xi'}{\xi}(\alpha)$.

\section{The clearest path to the Riemann Hypothesis}\label{ov1:sec:path}

The two volumes test many routes to RH and none of them closes here. Read together, they also show which route is clearest, what any proof along it must contain, and where it stops. We set this out here, before the details, so that each later chapter can be read against it.

\paragraph{The descent in level} For $\alpha>\frac12$ the shifted clock $\xi(\alpha)/\xi(\alpha+\sqrt u)$ is the Laplace transform of a law with no hypothesis on the zeros (Theorem~\ref{bp:thm:alpha}, which uses the verification of RH below height $1000$), and the density $v_\alpha/\pi$ is the Cauchy smoothing of the zeros at widths $\alpha-\beta$ (Proposition~\ref{prop:struct}); it is the Thorin density of that law whenever $\alpha$ lies to the right of every zero (Theorem~\ref{thm:poisson}). Two theorems turn RH into a single monotone question about this family.
\begin{enumerate}
\item Positivity is a zero free half plane (Theorem~\ref{cx:thm:positivity} of Volume II): for $\frac12\le\theta<1$, the density at level $\theta$ is a positive Thorin measure exactly when $\zeta$ has no zero with $\re s>\theta$. At $\theta=\frac12$ the equivalence with $\re(\xi'/\xi)>0$ on $\re s>\frac12$ is Lagarias's theorem.
\item The minimum sits at the origin (Theorem~\ref{thm:dbeta:minall}): RH holds if and only if $v_\alpha(t)>b_\alpha$ for every $\alpha\in(\frac12,1)$ and every $t\ne0$, and the same statement for $L(s,\chi_4)$ and $L(s,\chi_{12})$ is GRH for those characters.
\end{enumerate}
At the top of the strip the second statement is a theorem: $v_1(t)>v_1(0)$ for every $t\ne0$, with no hypothesis (Proposition~\ref{bp:prop:v1min}), so the clock at the basepoint one splits as a positive stable law of index $\frac12$ plus a GGC. For the two Dirichlet characters it holds for every $\alpha\ge\frac12+\frac1{\sqrt2}$. RH is the statement that the minimum at the origin, which holds at $\alpha=1$, persists as $\alpha$ descends to $\frac12$. In the language of Volume II, the Cauchy flow from the level one density, which the Euler product makes computable, can be run backwards to the critical line through positive measures (Theorem~\ref{tf:thm:deconv} of Volume II).

\paragraph{Three constraints on any proof} The volumes prove three facts that rule out whole classes of argument, and a route is worth pursuing only if it survives all three.
\begin{enumerate}
\item It cannot be finite. A perturbation of $\xi$ with the same zeros below a large height and four zeros off the line above it has an infinitely divisible, self decomposable reciprocal clock whose heat trace is strictly log convex and agrees with that of $\xi$ to the constant term (Theorem~\ref{thm:ch13:perturb} of Volume II); finite positivity tests on the heat trace, such as its leading Hankel blocks and its derivative signs to order $10^{17}$, pass for it. A computation can refute RH; it cannot prove it.
\item It must use the Euler product. The Davenport and Heilbronn function has a functional equation of the same shape, a positive trace and an infinitely divisible reciprocal law, and a zero off the line near $0.8085+85.6993\,i$ (Theorem~\ref{bsy:thm:dhid} of Volume II). Any argument that uses only the functional equation, positivity of a trace and infinite divisibility proves a false statement for it. Ki's approximation theorems point the same way (Chapter~\ref{chap:ki}).
\item It must not rest on a property that survives a Gaussian tilt. If a class of kernels is closed under $K\mapsto Ke^{-bu^2}$ and all its members have Fourier transforms with only real zeros, the P\'olya kernel is not in it (Proposition~\ref{ki:prop:obstruction}), because the de Bruijn Newman constant is not negative. Log concavity of the kernel and every other upper bound on $(\log K)''$ fall here.
\end{enumerate}
The descent in level survives the three. It is an analytic statement for all heights at once; its input at level one is the Euler product, whose prime sum fixes the density on the line $\re s=1$, and it fails for the Davenport and Heilbronn function; and it is a statement about where the zeros are, not a convexity property of the kernel.

\paragraph{Where it stops} The first step below the top is already the quasi Riemann hypothesis. Positivity at any level $\theta<1$ is a zero free half plane $\re s>\theta$ (statement 1 above), and no such half plane is known: the regions of Vinogradov and Korobov narrow towards the line $\re s=1$ as the height grows, and they are the best available. The density at level one is a smooth gamma part minus a compensated prime sum (Chapter~\ref{chap:lo} of Volume II); pushing its minimum downward is a statement about the size of that prime sum against the gamma part, which is, in Fourier variables, the error term $\psi(x)-x=O(x^{\theta+\varepsilon})$. Every other route in the two volumes, through the heat trace, Weil positivity in windows, hitting times, ferromagnetic approximations, truncations or friable approximants, either reduces to this step or meets one of the three constraints, and the chapters say which.

Two things would change the picture. If the zero free half plane $\re s>\frac78$ announced in October 2026 is confirmed, the Thorin property holds at every basepoint above $\frac78$ (Section~\ref{lv:sec:qrh}), and the open step becomes the descent from $\frac78$ to $\frac12$, a problem of the same kind and no easier in principle. And the de Bruijn Newman flow gives the only known analytic control below level one: at time $t=0.2$ every zero is real \cite{PT21}, and the Thorin phase diagram of Volume II (Chapter~\ref{chap:phase}) shows that what separates that theorem from RH is the existence of an analytic tail at positive time and the lack of a known one at time zero, not anything a computed density shows.

The clearest path is therefore to prove a zero free half plane by a positivity argument that uses the Euler product, and to carry it down to the line. It is honest to say that this is the hard core of the problem and not a way around it: the theorems of these volumes make the core precise, compute everything around it, and show that the remaining step is analytic. The penultimate chapter of this volume (Sections~\ref{route:sec:target} to~\ref{route:sec:answer}) assembles this route from the whole book, in six equivalent forms, compares it with curves over finite fields, where the same route is a proof, and states the missing step as a single positivity problem.
